%% ---------------------------------------------------------------------------
%% Ancres métriques éparses pour un a priori génératif 3D mono-vue
%% Version française du papier au format ICRA (IEEEtran, conference).
%%
%% ICRA se soumet en anglais : cette version sert au mémoire et à la relecture
%% interne. La version de soumission est icra_en.tex, et les deux doivent rester
%% synchronisées (mêmes chiffres, mêmes tableaux, même structure).
%%
%% Compilation :  tectonic -X compile icra_fr.tex
%% Piège : le caractère « œ » n'existe pas dans la police T1 utilisée et
%% disparaissait SILENCIEUSEMENT du PDF. Utiliser \oe{} et non le caractère brut.
%%
%% Recadré le 2026-09-11. La version antérieure (« Le toucher comme contrainte,
%% pas comme générateur ») mettait le toucher en premier et le problème du repère
%% en second. Les données disent l'inverse : le repère est le verrou de TOUTE ancre
%% métrique éparse, toucher et profondeur confondus, et il gouverne autant le gain
%% moyen que la variance.
%% ---------------------------------------------------------------------------
\documentclass[conference]{IEEEtran}
\IEEEoverridecommandlockouts

\usepackage[T1]{fontenc}
\usepackage[utf8]{inputenc}
\usepackage[french]{babel}
\usepackage{amsmath,amssymb}
\usepackage{graphicx}
\usepackage{booktabs}
\usepackage{multirow}
\usepackage{xcolor}
\usepackage{balance}
\usepackage[hidelinks]{hyperref}
\usepackage{cite}
\usepackage{tikz}
\usetikzlibrary{arrows.meta,positioning,fit,backgrounds,calc}
\usepackage{pgfplots}
\pgfplotsset{compat=1.18}

\definecolor{accent}{HTML}{2C5B79}
\definecolor{accentL}{HTML}{E7EFF5}
\definecolor{good}{HTML}{2A6B58}
\definecolor{goodL}{HTML}{E3EFEA}
\definecolor{bad}{HTML}{96412F}
\definecolor{lineC}{HTML}{95A1AD}

% --- libelles pour fig_method.tex ---
\newcommand{\lblUnet}{réseau de\\débruitage}
\newcommand{\lblCond}{image unique}
\newcommand{\lblDec}{décodage\\{\scriptsize ondelettes $\rightarrow$ SDF}}
\newcommand{\lblField}{champ de distance signée}
\newcommand{\lblContacts}{ancres $p_i$\\{\scriptsize mesurées}}
\newcommand{\lblInterp}{trilinéaire\\{\scriptsize écrite à la main}}
\newcommand{\lblGrad}{correction sur $\hat{z}_0$}
\newcommand{\lblLoop}{pas de débruitage suivant}
\newcommand{\lblMc}{marching\\cubes}
\newcommand{\lblNote}{une fois, à la fin\\(non différentiable)}
\newcommand{\lblBoxA}{différentiable, à chaque pas}
\newcommand{\lblBoxB}{aucune prédiction n'est mélangée à une autre}

% --- libelles pour fig_results.tex ---
\newcommand{\rlblZones}{4 zones au lieu de 8}
\newcommand{\rlblNoise}{bruit de placement 3\,mm / $15^\circ$}
\newcommand{\rlblPatchHi}{32 plaques au lieu de 8}
\newcommand{\rlblDens}{512 ancres au lieu de 32}
\newcommand{\rlblVis}{ancres hors de la face cachée}
\newcommand{\rlblNorm}{normales ajoutées}
\newcommand{\rlblPatchLo}{8 plaques au lieu de 4}
\newcommand{\rlblXlab}{écart apparié de F-Score@2\,\%}
\newcommand{\rlblPanelA}{Ce qu'un jeu d'ancres exige, et ce qu'il n'exige pas}
\newcommand{\rlblPanelB}{La visibilité ne prédit pas le gain}
\newcommand{\rlblXlabB}{part d'ancres occultées}
\newcommand{\rlblYlabB}{gain, F@2\,\%}
\newcommand{\rlblZero}{aucune ancre sur\\la face cachée}

\newcommand{\gd}[1]{\textbf{#1}}
\newcommand{\lo}[1]{\textbf{#1}}
\newcommand{\Fs}{F@2\,\%}

\title{\LARGE\bf Ancres métriques éparses pour un a priori génératif 3D mono-vue :\\
le verrou est le repère de sortie}

\author{%
  \IEEEauthorblockN{Nom de l'auteur}
  \IEEEauthorblockA{Institut des Systèmes Intelligents et de Robotique (ISIR)\\
  Sorbonne Université, Paris, France\\
  \texttt{auteur@isir.upmc.fr}}%
}

\begin{document}
\maketitle

\begin{abstract}
Un modèle génératif 3D mono-vue fournit un a priori de forme fort ; un robot fournit
quelques mesures métriques --- une poignée de contacts, une carte de profondeur. Nous
demandons ce qui limite l'injection de mesures aussi éparses dans un a priori
pré-entraîné, et trouvons que ce n'est ni leur nombre, ni leur contenu, ni la face de
l'objet qu'elles touchent. C'est le \emph{repère de sortie} du générateur. Sur un
modèle de diffusion à latent d'ondelettes d'un milliard de paramètres, nous montrons
que la pose de la forme générée n'est pas une fonction de l'image d'entrée : elle est
rééchantillonnée avec la graine de bruit, et se déplace de $130^\circ$ en médiane
d'une graine à l'autre sur 42 objets YCB. Deux voies d'ancrage sont ensuite mesurées
contre ce fait. Injecter des contacts comme contrainte différentiable pendant
l'échantillonnage, sans entraînement, améliore le F-Score@2\% de $+3{,}21$ ($p =
3{,}1\times10^{-4}$, test des signes), contre $+8{,}05$ lorsque le repère est donné :
l'estimation du repère consomme 60 à 74~\% du gain atteignable, et 91~\% de la
variance d'une exécution à l'autre. Ancrer le repère par conditionnement, en
reprojetant une seule carte de profondeur dans les vues canoniques d'une variante
multi-vues, fixe la pose à $\pm 2^\circ$ et gagne $+11{,}2$ en moyenne --- mais
l'effet est bimodal, perdant $-22$ sur les solides de révolution, et aucun sélecteur
sans vérité terrain testé n'atteint les $72{,}2$ qu'un oracle par objet obtient.
Empiler les contacts sur le repère ancré ajoute $+0{,}38$. Sur sept manipulations
contrôlées, ce qu'un jeu d'ancres exige est l'étalement spatial et la justesse du
placement ; le nombre, les normales et la visibilité ne comptent pas, et des ancres
sur la face que la caméra voit déjà valent autant que des ancres sur celle qu'elle ne
voit pas. Cinq formulations de fusion alternatives qui traitent les mesures comme un
générateur échouent toutes pour cette seule cause, que nous quantifions sans aucun
modèle de diffusion : six contacts orientés reconstruits par Poisson atteignent
$19{,}4$ \Fs{} contre $56{,}5$ pour l'image seule.
\end{abstract}

% ===========================================================================
\section{Introduction}
% ===========================================================================

Les grands modèles génératifs 3D pré-entraînés reconstruisent un objet plausible à
partir d'une image unique. Un robot qui tient cet objet sait des choses que l'image ne
dit pas : où sont ses doigts, et la profondeur de ce que la caméra voit. La question
naturelle est comment injecter ces mesures dans un modèle qui n'a jamais été entraîné
à les recevoir.

Ce travail a commencé comme une étude du toucher --- une poignée de contacts tactiles
peut-elle améliorer une reconstruction mono-image --- et son résultat central est que
la question était posée sur le mauvais axe. Les contacts aident, modestement et
significativement. Mais ce qui décide de combien ils aident n'a rien à voir avec les
contacts : ni leur nombre sur une plage de $16\times$, ni la présence d'une normale,
ni le fait qu'ils soient sur la face que la caméra ne voit pas. Ce qui décide est de
savoir si les mesures peuvent être placées dans le repère où le générateur produit sa
forme --- et ce repère, nous le montrons, n'est pas une fonction de l'entrée. Il est
tiré avec le bruit.

Nous traitons donc le problème comme un problème d'\emph{ancrage} d'un a priori par
des mesures métriques éparses, et nous utilisons deux sources de mesure comme deux
sondes du même verrou. Les contacts tactiles entrent comme contrainte différentiable
appliquée pendant l'échantillonnage, sans entraînement. Une carte de profondeur unique
entre comme signal de conditionnement, reprojetée dans les vues canoniques d'une
variante multi-vues de la même ossature. La première ne peut pas fixer le repère et le
paie ; la seconde le fixe, mais en tendant au modèle une preuve qu'il traite comme
réelle, avec des conséquences qui dépendent de l'objet. Aucune n'est une solution. Les
deux ensemble localisent le problème.

\subsection{Contributions}
\begin{enumerate}
  \item \textbf{La pose de sortie d'un modèle de diffusion 3D mono-vue n'est pas
  apprenable depuis son entrée} : sur trois graines de bruit et la même image, elle se
  déplace de $130^\circ$ en médiane sur 42 objets, les objets asymétriques n'étant pas
  plus stables que les symétriques (Sec.~\ref{sec:frame}).
  \item Un \textbf{schéma d'ancrage par contrainte sans entraînement}, différentiable
  de bout en bout à l'échantillonnage, et la mesure de ce que le repère lui coûte :
  $+3{,}21$ réalisé contre $+8{,}05$ repère donné
  (Secs.~\ref{sec:constraint}, \ref{sec:res-constraint}).
  \item Un \textbf{schéma d'ancrage par conditionnement} par reprojection de
  profondeur, qui fixe le repère et gagne $+11{,}2$, avec le constat que le gain est
  bimodal et que le sélecteur sans vérité terrain évident échoue
  (Secs.~\ref{sec:conditioning}, \ref{sec:res-conditioning}).
  \item Une \textbf{caractérisation de ce qu'un jeu d'ancres doit être} : sept
  manipulations contrôlées montrent que l'étalement et la justesse du placement
  comptent, tandis que le nombre, les normales et la visibilité ne comptent pas
  (Sec.~\ref{sec:res-anchor}).
  \item Une \textbf{décomposition de la variance} d'une exécution à l'autre, montrant
  que 91~\% en vient de la réestimation du repère et non du guidage, et trois pièges
  d'évaluation qui ont chacun renversé silencieusement une conclusion
  (Sec.~\ref{sec:methodology}).
\end{enumerate}

% ===========================================================================
\section{Travaux connexes}
% ===========================================================================

\textbf{Reconstruction de forme visuo-tactile.} Les travaux antérieurs combinent les
deux modalités soit en fusionnant des représentations apprises~\cite{smith2020,smith2021},
soit en accumulant des contacts denses dans une surface implicite~\cite{suresh2022}.
Les deux supposent assez de données tactiles pour décrire une forme :
\cite{smith2020} apprend une fusion par cartes à partir de touchers simulés sur un
grand ensemble d'objets, et \cite{suresh2022} intègre des relevés GelSight denses dans
une surface implicite par processus gaussien. Notre régime est l'inverse. Une saisie
avec une main à neuf zones rend un nombre de contacts à un seul chiffre, trop peu pour
ajuster une surface ; leur valeur est d'être des ancres sur un a priori, non un signal
à reconstruire. La Sec.~\ref{sec:fusions} montre que les formulations qui les traitent
en générateur échouent toutes pour une même raison.

\textbf{Modèles génératifs 3D.} Nous partons de WaLa~\cite{wala}, un modèle de
diffusion à latent d'ondelettes d'un milliard de paramètres qui prédit directement
$\hat{z}_0$ et le décode en champ de distance signée $256^3$. Sa variante mono-vue est
notre ossature principale ; sa variante à quatre vues de profondeur est le véhicule de
la Sec.~\ref{sec:conditioning}. Nous ne réentraînons rien.

\textbf{Guidage dans les modèles de diffusion.} Le guidage par classifieur
\cite{dhariwal} et le guidage sans classifieur~\cite{cfg} orientent l'échantillonnage
vers une condition. Notre schéma par contrainte est le plus proche du premier par
l'esprit, mais n'utilise aucun classifieur : la correction est le gradient analytique
d'un résidu géométrique, donc sans paramètre appris ni distribution d'entraînement à
quitter.

\textbf{Génération conditionnée par la pose.} Zero-1-to-3~\cite{zero123} conditionne
sur une transformation caméra relative explicite ; les méthodes stéréo par passe avant
comme DUSt3R~\cite{dust3r} régressent la pose directement depuis une paire d'images.
Une initialisation façon SDEdit~\cite{sdedit} offre une troisième voie, que nous
discutons sans la tester. La Sec.~\ref{sec:frame} explique pourquoi aucune n'est un
correctif immédiat ici : l'ossature mono-vue ne reçoit aucune information de caméra, et
les variantes qui en reçoivent détruisent précisément la marge que les mesures devaient
combler.

% ===========================================================================
\section{Le repère de sortie n'est pas une fonction de l'entrée}
\label{sec:frame}
% ===========================================================================

Les mesures vivent dans le repère du capteur ; le générateur produit dans un repère
canonique qui lui est propre. Tout schéma d'ancrage doit envoyer l'un sur l'autre.
Avant de demander avec quelle précision, nous demandons si la cible de cette
application existe.

\textbf{La cause, en une ligne de configuration.} L'ossature mono-vue déclare
\texttt{use\_camera\_index = False} : la fonction qui extrait l'indice de vue renvoie
\texttt{None}, et le modèle ne reçoit aucune information sur la caméra. Il n'y a donc
aucune convention de repère à retrouver ; la pose de sortie est ce que le contenu de
l'image et le bruit produisent.

\textbf{La pose est rééchantillonnée avec le bruit.} Le correctif le moins coûteux
imaginable serait de régresser la rotation de sortie depuis l'image, sans réentraîner
le générateur. Il présuppose que la pose soit une fonction de l'image. En comparant les
rotations obtenues pour trois graines de bruit sur la \emph{même} image, sur 42
objets, l'écart médian entre graines vaut $140{,}5^\circ$ pour les objets asymétriques
et $129{,}8^\circ$ pour les symétriques ; 1 objet asymétrique sur 16 est stable sous
$30^\circ$, 0 sur 26 chez les symétriques ; la corrélation asymétrie / stabilité vaut
$-0{,}073$. Le partage par symétrie écarte un artefact de mesure --- deux sorties
identiques d'un solide de révolution se recalent avec des rotations très différentes
--- et les objets asymétriques, dont l'orientation est identifiable, ne sont \emph{pas
plus stables}. Aucun régresseur ne peut apprendre cette pose : non faute du bon
estimateur, mais parce que la cible n'existe pas.

\textbf{Les modèles qui reçoivent la pose l'ancrent.} Les variantes multi-vues de la
même ossature déclarent \texttt{use\_camera\_index = True} et prennent quatre images
avec leur indice dans un rig canonique à 55 positions. Leur rotation entre prédiction
et référence est reproductible à $0{,}8^\circ$ d'une graine à l'autre et identique sur
tous les objets bien contraints. Mais quatre vues résolvent aussi la reconstruction
--- jusqu'à $99{,}9$ \Fs{} sur des objets que le mono-vue note entre $23$ et $59$ ---
ne laissant rien à ajouter à une mesure. \textbf{Pose et marge ne coexistent dans
aucune variante publiée.} C'est la tension que les deux voies d'ancrage ci-dessous
doivent négocier.

% ===========================================================================
\section{Deux façons d'ancrer}
% ===========================================================================

\subsection{Ancrage par contrainte pendant l'échantillonnage}
\label{sec:constraint}

L'ossature prédit directement $\hat{z}_0$ plutôt que le résidu de bruit : la sortie du
réseau \emph{est} donc l'estimation de forme courante, et peut être corrigée sans
toucher au calendrier de bruit ni à l'échantillonneur. À chaque pas, nous décodons
$\hat{z}_0$ en champ de distance signée dense $\Phi(\hat{z}_0) \in
\mathbb{R}^{256^3}$, l'évaluons aux $N$ points d'ancrage $p_i$, et avançons selon le
gradient de

\begin{equation}
  \mathcal{L}_{\text{pos}} \;=\; \frac{1}{N}\sum_{i=1}^{N}
  \Big[\, \mathcal{T}\big(\Phi(\hat{z}_0),\, p_i \big) \Big]^2 ,
  \label{eq:loss}
\end{equation}

où $\mathcal{T}$ est l'interpolation trilinéaire. La contrainte imposée est simplement
\emph{la surface passe par ces points} (Fig.~\ref{fig:method}). Le pas est normalisé,
$\Delta = w \,\lVert \hat{z}_0 \rVert \, \nabla \mathcal{L} / \lVert \nabla
\mathcal{L}\rVert$, pour que le poids $w$ ne dépende ni de l'échelle du champ ni de
celle du latent, qui dérivent tous deux pendant l'échantillonnage.

\begin{figure*}[t]
\centering
\resizebox{\textwidth}{!}{%% ---------------------------------------------------------------------------
%% fig_method.tex -- method figure, shared geometry between icra_en and icra_fr.
%%
%% The calling document must define the label macros BEFORE \input-ing this file:
%%   \lblUnet \lblCond \lblDec \lblField \lblContacts \lblInterp \lblGrad
%%   \lblLoop \lblMc \lblNote \lblBoxA \lblBoxB
%% See the definitions in either paper for the expected wording.
%%
%% One figure, one claim: the correction acts on zhat_0 INSIDE the loop, through
%% a differentiable decode. Marching cubes sits outside, because it is not
%% differentiable -- and the guidance never needs it.
%%
%% Layout: two rows, columns aligned so every arrow is orthogonal. The two rows
%% are joined on the right (field -> interpolation) and on the left (correction
%% -> zhat_0), which is what makes the loop legible.
%% ---------------------------------------------------------------------------

\begin{tikzpicture}[
  font=\small,
  >={Stealth[length=2.4mm]},
  box/.style    ={draw=lineC, fill=white, rounded corners=1.5pt, inner sep=5pt,
                  align=center, minimum height=9mm},
  op/.style     ={box, fill=accentL, draw=accent},
  data/.style   ={box, fill=white, draw=lineC},
  meas/.style   ={box, fill=goodL, draw=good, very thick},
  export/.style ={box, fill=black!5, draw=lineC, dashed},
  flow/.style   ={->, draw=accent, thick},
  back/.style   ={->, draw=bad, very thick},
  side/.style   ={->, draw=good, very thick},
  lbl/.style    ={font=\scriptsize, inner sep=2pt},
]

% ---- geometry --------------------------------------------------------------
\def\rowA{0}        % sampler row
\def\rowB{-2.7}     % geometric-residual row
\def\cA{0}   \def\cB{2.3}  \def\cC{4.9}
\def\cD{8.3} \def\cE{11.3} \def\cF{14.3}

% ---- top row: the sampler --------------------------------------------------
\node[data] (xt)   at (\cA,\rowA) {$x_t$};
\node[op]   (unet) at (\cB,\rowA) {\lblUnet};
\node[data] (z)    at (\cC,\rowA) {$\hat{z}_0$};
\node[op]   (dec)  at (\cD,\rowA) {\lblDec};
\node[data] (phi)  at (\cE,\rowA) {$\Phi \in \mathbb{R}^{256^3}$\\[-2pt]
                                   {\scriptsize\lblField}};

\draw[flow] (xt)   -- (unet);
\draw[flow] (unet) -- (z);
\draw[flow] (z)    -- (dec);
\draw[flow] (dec)  -- (phi);

% conditioning image enters the network from below, clear of the loop arrow
\node[lbl, align=center] (cond) at (\cB,-1.35) {\lblCond};
\draw[flow] (cond) -- (unet);

% ---- bottom row: the geometric residual ------------------------------------
\node[meas] (pts)  at (\cF,\rowB) {\lblContacts};
\node[op]   (int)  at (\cE,\rowB) {\lblInterp};
\node[op]   (loss) at (\cD,\rowB) {$\mathcal{L}=\frac{1}{N}\sum_i \Phi(p_i)^2$};
\node[op]   (grad) at (\cC,\rowB) {$\hat{z}_0 - w\,\lVert\hat{z}_0\rVert\,\hat{g}$\\[-1pt]
                                   {\scriptsize $\hat{g}=\nabla\mathcal{L}/\lVert\nabla\mathcal{L}\rVert$}};

\draw[side] (pts)  -- (int);
\draw[flow] (phi)  -- (int);          % vertical: columns are aligned
\draw[flow] (int)  -- (loss);
\draw[back] (loss) -- (grad);

% the correction returns to zhat_0 -- the whole point of the figure
\draw[back] (grad) -- (z)
  node[lbl, midway, right, text=bad, align=left] {\lblGrad};

% ---- the loop --------------------------------------------------------------
\draw[flow] (z.north) -- ++(0,0.95) -| (xt.north)
  node[lbl, pos=0.25, above] {\lblLoop};

% ---- outside the loop: mesh export -----------------------------------------
\node[export] (mc) at (\cF,\rowA) {\lblMc};
\draw[flow, dashed] (phi) -- (mc);
\node[lbl, above=1.5mm of mc, text=black!55, align=center] {\lblNote};

% ---- framing ---------------------------------------------------------------
\begin{scope}[on background layer]
  \node[fit=(grad)(loss)(int)(pts), draw=accent, dashed, rounded corners=3pt,
        inner sep=8pt, fill=accent!4] (boxA) {};
\end{scope}
\node[lbl, text=accent, anchor=south east, yshift=1mm] at (boxA.north east)
  {\lblBoxA};
\node[lbl, text=black!55, anchor=north east, yshift=-1mm] at (boxA.south east)
  {\lblBoxB};

\end{tikzpicture}
}
\caption{\textbf{L'ancrage par contrainte, à sa place.} La correction agit sur le
$\hat{z}_0$ prédit \emph{à l'intérieur} de la boucle de débruitage, à travers un
décodage différentiable de bout en bout ; les ancres sont la seule grandeur mesurée
du schéma (en vert). Les marching cubes sont hors de la boucle parce qu'ils ne sont pas
différentiables --- et le guidage n'a jamais besoin du maillage, seulement du champ.
Aucune prédiction n'étant jamais mélangée à une autre, l'interférence destructive qui
condamne les formulations par composition de la Sec.~\ref{sec:fusions} ne peut pas se
produire.}
\label{fig:method}
\end{figure*}

Le guidage n'a jamais besoin du maillage, seulement du champ : les marching
cubes~\cite{marchingcubes} ne sont appliqués qu'une fois, à la fin, pour exporter le
résultat. Point décisif : \emph{aucune prédiction n'est jamais mélangée à une autre}.
L'interférence destructive qui condamne les formulations par composition de la
Sec.~\ref{sec:fusions} ne peut pas se produire.

\textbf{Mise en \oe{}uvre.} $\mathcal{T}$ est écrite à la main, en indexation et
produits, plutôt qu'avec \texttt{grid\_sample} : le noyau rétropropagé 3D de cette
dernière n'est pas implémenté sur Apple Metal, et le repli CPU imposerait un
aller-retour de 67\,Mo par pas. Coût mesuré : $3{,}44$\,s par pas guidé, $8{,}98$\,Go
de crête mémoire, $\approx 55$\,s par échantillon à 20 pas sur un M2 Max.

\textbf{Terme optionnel de normale.} Un contact venant d'une main porte une normale
sortante approchée $n_i$. Nous n'imposons que la \emph{direction} du gradient du champ,
$\mathcal{L}_{\text{nrm}} = \frac{1}{N}\sum_i [\,1 - s\, g_i\!\cdot n_i / \lVert g_i
\rVert\,]$ avec $g_i = \nabla_{\!p}\Phi|_{p_i}$ par différences finies centrées et
$s = +1$ la convention de signe mesurée (extérieur positif : médiane intérieure
$-0{,}0998$, extérieure $+0{,}0999$, une TSDF tronquée à $\pm 0{,}10$). Les deux
pertes ont des unités incommensurables ; $\lambda$ est calibré une fois au premier pas
guidé pour que leurs normes de gradient coïncident, puis figé. La
Sec.~\ref{sec:res-anchor} montre que le terme n'aide pas de façon mesurable.

\textbf{Placer les ancres.} Sans vérité terrain, le repère est estimé par deux
critères : la \emph{cohérence de silhouette} (projeter la prédiction et la comparer à
l'image de conditionnement, que le modèle a déjà reçue) et le \emph{résidu des ancres}
(parmi les poses que la silhouette laisse plausibles, garder celle dont les ancres
s'approchent le plus de la surface de la première passe). Aucun ne suffit seul.
Recaler les ancres sur une première passe non guidée par ICP les place, comparées à
leur position exacte, à $0{,}66$--$0{,}76$ en médiane dans un repère où l'objet occupe
$[-1,1]$ --- \emph{sur} la surface mais \emph{au mauvais endroit}, l'ICP superposant
sans respecter les correspondances et les quasi-symétries lui en laissant la latitude.
Silhouette plus résidu ramènent l'écart médian de $0{,}755$ à $0{,}501$, et à
$0{,}111$ sur les objets sans ambiguïté de silhouette.

\subsection{Ancrage par conditionnement à travers une variante multi-vues}
\label{sec:conditioning}

La variante à quatre vues de profondeur ancre le repère mais exige quatre entrées, et
quatre vues réelles ne laissent aucune marge. Fournir une vue réelle et trois images
vides échoue : mesuré sur quatre objets et deux graines, $43{,}07$ \Fs{} contre
$95{,}40$ avec quatre vues réelles, et moins bien que le modèle mono-vue dédié sur la
moitié des objets. \textbf{Les vues vides ne sont pas neutres ; elles retirent de
l'information.}

Une carte de profondeur, en revanche, est de la géométrie 3D dans le repère de la
caméra. On peut la rétroprojeter en points, puis la revoir depuis n'importe quelle
autre position du rig canonique. La vue synthétique est géométriquement juste là où
elle a de la donnée et vide ailleurs --- bien plus proche de la distribution
d'entraînement qu'une image noire. Nous fournissons donc au modèle \emph{une vue de
profondeur réelle et trois reprojetées depuis elle}, avec leurs quatre indices de
caméra. Le placement des ancres devient alors une transformation fixe, calibrée une
fois : mesuré sur quatre objets et deux graines, la rotation de sortie reste entre
$118^\circ$ et $122^\circ$ autour d'un axe constant.

L'intention de conception était que la reprojection \emph{n'ajoute aucune
information} --- la même observation unique réencodée --- pour que la reconstruction
reste au niveau mono-vue pendant que les indices de caméra ancrent la pose. La
Sec.~\ref{sec:res-conditioning} montre que cette intention n'a pas été atteinte : le
modèle traite les vues reprojetées comme des preuves, et ce qu'il fait de leurs trous
dépend de l'objet.

\subsection{Modèle d'ancres}
\label{sec:contactmodel}

Échantillonner les ancres uniformément sur le maillage de référence est l'objection la
plus lourde contre tout résultat d'ancrage, puisque les points viennent alors de la
réponse. Nous le remplaçons par un modèle de palpation imposant trois contraintes
physiques : accessibilité en ligne droite (rayons lancés depuis l'extérieur, premiers
impacts seuls), restriction à la face occultée par la caméra, et rejet des points où
une pulpe de 6\,mm ne passerait pas. Les ancres sont regroupées en \emph{sites} plutôt
que dispersées --- une main touche par plaques. Ce seul changement divise par deux le
plafond de la méthode ($+14{,}37 \rightarrow +6{,}77$ sur cinq objets) : une partie du
gain annoncé par les campagnes intermédiaires venait de la couverture uniforme, pas des
mesures.

Pour la Sec.~\ref{sec:res-hand}, nous utilisons en outre un \emph{simulateur de
préhension} qui place une ancre par zone de doigt (pouce, index, majeur, annulaire,
auriculaire, paume) à des angles anatomiques fixes autour de l'axe principal de
l'objet, bien plus proche d'une main réelle que des sites bien répartis : écart minimal
entre ancres de $0{,}058$ à $0{,}344$ contre $0{,}446$ à $1{,}281$ pour nos sites.

% ===========================================================================
\section{Protocole expérimental}
% ===========================================================================

\textbf{Protocole.} 42 objets YCB~\cite{ycb}, 3 graines de bruit chacun, soit 126
comparaisons \emph{appariées} : chaque génération ancrée est comparée à une génération
non ancrée de \emph{même graine}. Rien n'étant entraîné, tous les objets sont des
objets de test.

\textbf{Métrique.} F-Score@2\% après recalage sur la référence, distance de Chamfer
symétrique, connectivité, et rappel séparé selon la visibilité caméra.

\textbf{Significativité.} Nous utilisons un \textbf{test des signes}, pas un test de
Student : la distribution des gains est fortement asymétrique (médiane $+1{,}69$ pour
une moyenne de $+3{,}21$) et quelques objets dominent. Une moyenne y est fragile, un
compte de signes ne l'est pas. Tous les $p$ sont \textbf{bilatéraux}, et les ex aequo
exacts sont écartés plutôt que comptés comme des échecs --- un cas sur 126 est noté à
l'identique avec et sans ancrage, d'où un résultat principal rapporté sur 125.

\textbf{Discipline de l'oracle.} Là où c'est indiqué, les ancres sont placées à l'aide
du maillage de référence, tirées avec le modèle de palpation de la Sec.~\ref{sec:contactmodel}. Ces chiffres sont rapportés \emph{uniquement} comme bornes
supérieures, jamais comme performance de système, et toujours avec mention explicite
de ce qui a été donné gratuitement.

% ===========================================================================
\section{Résultats}
% ===========================================================================

\subsection{Ancrage par contrainte : ce que le repère coûte}
\label{sec:res-constraint}

\begin{table}[t]
\centering
\caption{Ancrage par contrainte avec des contacts tactiles, 42 objets YCB.}
\label{tab:main}
\begin{tabular}{lr}
\toprule
Témoin non ancré (image seule) & 53,1 \Fs \\
\midrule
Gain moyen, repère estimé & \gd{$+3{,}21$} \\
Médiane / écart-type & $+1{,}69$ / $7{,}52$ \\
Cas positifs & \gd{83 / 125} \\
Test des signes (bilatéral) & \gd{$p = 3{,}1\times10^{-4}$} \\
Objets à gain moyen positif & 35 / 42 \\
\midrule
Gain moyen, repère donné (oracle) & $+8{,}05$ \\
\bottomrule
\end{tabular}
\end{table}

Le tableau~\ref{tab:main} donne le résultat par contrainte. L'effet est modeste au
regard de sa dispersion : c'est la significativité du \emph{signe} qui est solide, pas
la précision de l'amplitude. Sur cinq objets, une campagne antérieure mesurait
$+1{,}79$ à $p = 0{,}035$ ; à 42 objets l'effet est plus fort et deux ordres de
grandeur plus significatif --- ce n'était donc pas un objet qui portait tout.

Comparer le gain réalisé à l'oracle isole le coût de l'estimation du repère :
$+3{,}21$ contre $+8{,}05$ sur 42 objets, et $+1{,}79$ contre $+6{,}77$ sous le modèle
de palpation sur cinq. \textbf{L'estimation du repère absorbe 60 à 74~\% de ce que
l'ancrage par contrainte peut rendre} --- un ordre de grandeur au-dessus de tout autre
facteur mesuré dans ce papier.

\subsection{Ancrage par conditionnement : le repère fixé, à un prix}
\label{sec:res-conditioning}

\begin{table}[t]
\centering
\caption{Les deux voies d'ancrage, et leur combinaison. Mêmes 42 objets, mêmes 126
cas appariés par graine, même métrique.}
\label{tab:routes}
\begin{tabular}{lrr}
\toprule
& \Fs & repère \\
\midrule
Image seule (ossature mono-vue) & 53,1 & tiré avec le bruit \\
\quad $+$ contacts en contrainte & 56,3 & estimé, $\pm$ large \\
Profondeur reprojetée en 4 vues & \gd{64,3} & fixé, $\pm 2^\circ$ \\
\quad $+$ contacts en contrainte & 64,7 & fixé \\
\midrule
Choix oracle de la voie, par cas & 72,2 & --- \\
Quatre vues de profondeur \emph{réelles} (4 objets) & 98,6 & fixé \\
\bottomrule
\end{tabular}
\end{table}

Le tableau~\ref{tab:routes} place les deux voies côte à côte. Reprojeter une seule
carte de profondeur dans les quatre vues canoniques gagne $+11{,}22$ sur l'ossature
mono-vue en moyenne appariée --- plus de trois fois ce que les contacts rendent --- et
fixe le repère. Mais le gain est \textbf{bimodal} : 87 cas sur 126 s'améliorent de
$+26{,}1$ en moyenne, et 39 cas perdent $-22{,}1$. Les pertes sont catastrophiques sur
les solides de révolution --- balle de baseball $95{,}5 \rightarrow 29{,}2$, ballon
$93{,}9 \rightarrow 35{,}5$, conserve $91{,}0 \rightarrow 46{,}8$ --- où une seule vue
de profondeur laisse trois vues canoniques presque vides, que le modèle remplit depuis
son a priori, mal. L'intention de la Sec.~\ref{sec:conditioning}, que la reprojection
n'ajoute aucune information, n'a pas été atteinte : le modèle traite les vues
reprojetées comme des preuves.

\textbf{Aucun sélecteur sans vérité terrain testé n'atteint l'oracle.} Un choix par
cas entre les deux voies atteindrait $72{,}2$. Le sélecteur évident --- la part de
pixels reprojetés qui portent une profondeur --- est \emph{anti}-corrélé au gain
($r = -0{,}40$) : les objets compacts et ronds remplissent beaucoup de pixels et
échouent, les objets plats en remplissent peu et réussissent. Le remplissage mesure la
forme de l'objet, pas la qualité de la reprojection, et aucun seuil ne bat le choix
systématique de la reprojection. La classification par symétrie, l'autre candidat, est
déjà réfutée comme prédicteur en Sec.~\ref{sec:res-anchor}. Le $72{,}2$ existe mais
est hors de portée sans connaître la réponse.

\textbf{Empiler les contacts sur le repère ancré ajoute $+0{,}38$} (75/126 positifs,
$p = 0{,}040$). Une fois le repère fixé et l'a priori fort, les contacts n'ont plus
grand-chose à corriger. Les deux voies ne sont pas complémentaires non plus : la
corrélation entre le gain de reprojection et le gain de contacts vaut $+0{,}056$, et
sur les 39 cas où la reprojection échoue, les contacts récupèrent $+1{,}7$ contre une
perte de $-22$. Elles échouent sur des objets différents par accident, pas par
conception.

\subsection{Ce qu'un jeu d'ancres exige}
\label{sec:res-anchor}

\begin{figure*}[t]
\centering
\resizebox{\textwidth}{!}{%% ---------------------------------------------------------------------------
%% fig_results.tex -- results figure, shared geometry between icra_en/icra_fr.
%%
%% Requires the label macros \rlbl* to be defined by the calling document.
%%
%% Two panels, one claim each, and they answer each other:
%%   (a) seven controlled manipulations, all paired. Two families of standard
%%       error: 0.58 within a campaign (pooled over four objects; shared oracle frame, placement noise
%%       cancels) and 0.91 across campaigns (no shared unanchored pass). Colour
%%       encodes the two-sided sign test, which carries the conclusions. Three
%%       clear it, four do not -- and the three that do are all about WHERE the
%%       anchors are and how accurately they are placed, never how many there
%%       are or what the camera can see of them.
%%   (b) the visibility result on its own, because it is the one that runs
%%       against the premise of the work. Visibility is decided from the camera
%%       that rendered the images (tools/visibility_render_camera.py); the
%%       campaign's own fractions came from a camera placed in the wrong frame.
%%
%% All values are paired differences between two guided conditions on the same
%% object and the same seed; n = 15 each.
%% ---------------------------------------------------------------------------

\begin{tikzpicture}

% ===== panel (a) : forest plot ==============================================
\begin{axis}[
  name=fa,
  width=9.2cm, height=5.4cm,
  scale only axis,
  xmin=-4.2, xmax=4.9,
  ymin=-0.8, ymax=6.8,
  ytick={0,1,2,3,4,5,6},
  yticklabels={\rlblZones,\rlblNoise,\rlblPatchHi,\rlblDens,\rlblVis,\rlblNorm,\rlblPatchLo},
  yticklabel style={font=\scriptsize, align=right},
  xtick={-4,-2,0,2},
  xticklabel style={font=\scriptsize},
  xlabel={\scriptsize \rlblXlab},
  xlabel style={yshift=2pt},
  axis line style={draw=black!45},
  tick style={draw=black!45},
  grid=none,
  clip=false,
]
  \draw[black!45] (axis cs:0,-0.8) -- (axis cs:0,6.8);

  % Barres d'erreur a une erreur type. ATTENTION, deux familles :
  %   0,58 quand les deux conditions partagent le repere oracle (meme campagne :
  %        le bruit de placement s'annule dans la difference appariee) ;
  %   0,91 quand elles viennent de campagnes distinctes, qui ne partagent aucune
  %        passe non guidee, et ou ce bruit ne s'annule pas.
  % La couleur encode le test des signes bilateral, pas ces erreurs types : c'est
  % lui qui porte les conclusions, parce qu'il ne demande aucune estimation de sigma.
  \addplot[only marks, mark=*, mark size=1.7pt, draw=accent, fill=accent,
           error bars/.cd, x dir=both, x explicit, error bar style={accent}]
    table[x=v, y=i, x error=e] {
      v      i   e
      -2.55  0   0.58
      -2.40  1   0.58
       2.22  6   0.91
    };
  \addplot[only marks, mark=*, mark size=1.7pt, draw=black!45, fill=black!30,
           error bars/.cd, x dir=both, x explicit, error bar style={black!45}]
    table[x=v, y=i, x error=e] {
      v      i   e
      -0.24  2   0.91
      -0.16  3   0.91
       0.82  4   0.58
       0.93  5   0.58
    };

  % le compte de signes, qui est l'evidence reelle
  \node[font=\scriptsize, text=accent,   anchor=west] at (axis cs:3.0,0) {3/15};
  \node[font=\scriptsize, text=accent,   anchor=west] at (axis cs:3.0,1) {3/15};
  \node[font=\scriptsize, text=black!55, anchor=west] at (axis cs:3.0,2) {5/15};
  \node[font=\scriptsize, text=black!55, anchor=west] at (axis cs:3.0,3) {8/15};
  \node[font=\scriptsize, text=black!55, anchor=west] at (axis cs:3.0,4) {11/15};
  \node[font=\scriptsize, text=black!55, anchor=west] at (axis cs:3.0,5) {10/15};
  \node[font=\scriptsize, text=accent,   anchor=west] at (axis cs:3.0,6) {12/15};

\end{axis}

\node[font=\scriptsize\bfseries, anchor=south west] at (fa.north west)
  {(a)\; \rlblPanelA};

% ===== panel (b) : occlusion vs gain ========================================
\begin{axis}[
  name=fb,
  at={(fa.right of east)}, anchor=left of west, xshift=14mm,
  width=5.4cm, height=5.4cm,
  scale only axis,
  xmin=0.08, xmax=1.08,
  ymin=-12, ymax=13,
  xtick={0.25,0.5,0.75,1},
  xticklabel style={font=\scriptsize},
  ytick={-10,-5,0,5,10},
  yticklabel style={font=\scriptsize},
  xlabel={\scriptsize \rlblXlabB},
  ylabel={\scriptsize \rlblYlabB},
  xlabel style={yshift=2pt}, ylabel style={yshift=-4pt},
  axis line style={draw=black!45},
  tick style={draw=black!45},
  clip=false,
]
  \draw[black!30] (axis cs:0.08,0) -- (axis cs:1.08,0);

  % the 45 cases, hidden fraction seen from the rendering camera
  \addplot[only marks, mark=o, mark size=1.25pt, draw=accent!75] coordinates {
    (0.333,1.45) (0.333,0.68) (0.333,0.75) (0.833,0.97) (0.833,0.84) (0.833,0.57)
    (0.167,1.64) (0.167,3.29) (0.167,1.01) (0.333,3.59) (0.333,-1.74) (0.333,3.67)
    (0.833,3.31) (0.833,2.29) (0.833,0.09) (0.167,4.18) (0.167,3.11) (0.167,8.15)
    (0.833,0.42) (0.833,-1.04) (0.833,0.67) (0.833,-10.43) (0.833,-1.25) (0.833,-0.92)
    (0.167,-3.98) (0.167,-0.72) (0.167,-0.87) (0.667,5.66) (0.667,-1.14) (0.667,-6.93)
    (0.833,7.02) (0.833,-5.67) (0.833,5.88) (0.667,7.60) (0.667,0.13) (0.667,-2.27)
    (0.833,10.12) (0.833,0.01) (0.833,9.76) (1.000,11.55) (1.000,0.94) (1.000,5.38)
    (0.333,9.16) (0.333,0.52) (0.333,7.36)
  };
  % least-squares fit: flat
  \addplot[draw=bad, very thick, domain=0.15:1.0, samples=2] {-0.013*x + 1.893};

  \node[font=\scriptsize, text=bad, anchor=north west] at (axis cs:0.10,12.4)
    {$r=-0.001$};
\end{axis}

\node[font=\scriptsize\bfseries, anchor=south west] at (fb.north west)
  {(b)\; \rlblPanelB};

\end{tikzpicture}
}
\caption{\textbf{Sept manipulations contrôlées du jeu d'ancres.}
\textbf{(a)} Chaque estimation est une différence appariée entre deux conditions
ancrées, sur le même objet et la même graine de bruit ($n=15$). Les barres d'erreur
valent une erreur type --- $0{,}26$ lorsque les deux conditions viennent de la même
campagne et partagent donc le repère oracle, le bruit de placement s'annulant alors
dans la différence, et $0{,}91$ lorsqu'elles viennent de campagnes distinctes, qui ne
partagent aucune passe non ancrée. La couleur marque la significativité au test des
signes bilatéral, dont les comptes figurent à droite ; nous nous appuyons sur lui
plutôt que sur les erreurs types, car il ne demande aucune estimation de $\sigma$.
Trois comparaisons seulement en sortent, et toutes trois portent sur \emph{où} sont les
ancres et \emph{avec quelle justesse} elles sont placées --- jamais sur leur nombre, ni
sur ce que la caméra en voit. \textbf{(b)} Le résultat de visibilité seul, sur les 45
cas de la campagne de stratégies de préhension. L'ajustement par moindres carrés penche
dans le mauvais sens, et les trois cas où la main ne touche \emph{rien} que la caméra
ne voie déjà (en vert) comptent parmi les plus gros gains mesurés.}
\label{fig:results}
\end{figure*}

La Fig.~\ref{fig:results}a balaie sept propriétés du jeu d'ancres sous ancrage par
contrainte, sur cinq objets, placement oracle. Trois passent un test des signes
bilatéral à $p = 0{,}035$, et elles partagent un thème.

\textbf{L'étalement compte.} À budget constant de 128 ancres réparties sur $N$ sites,
le gain vaut $-0{,}42$ à 4 sites, $+1{,}79$ à 8, $+2{,}03$ à 16 et $+1{,}55$ à 32 : un
seuil entre 4 et 8, puis un plateau. Passer de 8 plaques à 4 coûte $-2{,}55$ (12/15
négatifs).

\textbf{La justesse du placement compte.} Un bruit gaussien de $3$\,mm en position et
$15^\circ$ sur les normales, l'ordre de ce que la cinématique directe d'un doigt
fournirait, coûte $-2{,}40$ (12/15 négatifs).

\textbf{Le nombre ne compte pas.} De 32 à 512 ancres, le gain vaut $+0{,}23$,
$+1{,}21$, $+0{,}07$ : plat sur une plage de $16\times$.

\textbf{Les normales ne comptent pas.} Les ajouter donne $+0{,}93$, que le test des
signes ne sépare pas de zéro ($10/15$, $p = 0{,}302$). La moyenne est portée par deux
valeurs aberrantes : retirer la plus grande laisse $+0{,}43$, pour une médiane de
$+0{,}26$.

\textbf{La visibilité ne compte pas --- et c'est le résultat qui renverse la
prémisse.} Avec le simulateur de préhension, trois stratégies placent les mêmes six
ancres avec $83$~\%, $70$~\% et $37$~\% d'entre elles sur la face occultée. Leurs gains
valent $+1{,}73$, $+1{,}37$ et $+2{,}55$ : la stratégie qui touche le \emph{moins} la
face cachée gagne le \emph{plus}. Sur 45 cas, la corrélation entre part occultée et
gain vaut $r = -0{,}145$, légèrement dans le mauvais sens (Fig.~\ref{fig:results}b).
Le cas décisif est sans ambiguïté : sur la perceuse, la stratégie qui place
\textbf{zéro} ancre sur la face occultée --- ne touchant que ce que la caméra voit déjà
--- gagne $+5{,}68$ en moyenne. Les ancres éparses ne complètent pas ce que la vision
manque. \textbf{Elles fixent position, échelle et orientation globales d'un a priori
dont le repère est lâche, et pour cela le côté où elles se trouvent n'importe pas.}

\subsection{Pourquoi les formulations de fusion échouent}
\label{sec:fusions}

Cinq formulations exigent que les mesures produisent une forme complète, ensuite
mélangée à celle de l'image. L'échec le plus instructif est le mélange pondéré dans
l'espace $\hat{z}_0$, qui s'effondre de $44{,}9$ à $13{,}9$ \Fs{} à poids égal :
l'interpolation convexe entre deux modes distincts d'une distribution tombe dans une
zone de faible densité --- le point médian entre deux formes plausibles n'est pas une
forme plausible. Cela vaut pour tout mélange naïf de latents génératifs.

\begin{table}[t]
\centering
\caption{Six contacts orientés, avec et sans a priori génératif.}
\label{tab:poisson}
\begin{tabular}{lr}
\toprule
Reconstruction de Poisson~\cite{poisson}, 6 contacts orientés & 19,41 \Fs \\
Reconstruction de Poisson, 18 contacts (3 prises) & 20,10 \\
\midrule
A priori génératif, image seule, sans mesure & \gd{56,52} \\
\bottomrule
\end{tabular}
\end{table}

Le tableau~\ref{tab:poisson} convertit la cause commune en une mesure, sans aucun
modèle de diffusion. La reconstruction de Poisson consomme précisément des points
orientés et rend une surface fermée dans 9 cas sur 10 --- ce n'est pas du gravier.
Elle n'atteint pourtant que 34~\% de ce que l'image seule obtient déjà. Les mesures
éparses sont un mauvais générateur et une bonne ancre ; la répartition des rôles n'est
pas un choix de conception mais ce que les données imposent.

Un module de fusion \emph{appris} (encodeur de géométrie et cross-attention greffés
sur l'ossature gelée, 50 époques, 5 à 10\,M de paramètres entraînés) mesure $-4{,}20$
sur cinq objets en tirages non appariés --- $0{,}9$ erreur type, indiscernable de
zéro. La comparaison avec l'ancrage par contrainte ne porte pas vraiment sur les deux
scores, tous deux médiocres. Elle porte sur le fait que la voie apprise
\textbf{consomme les objets dont elle aurait besoin pour se juger} : 15 sur 20
deviennent des données d'entraînement, il en reste cinq pour conclure, avec un bruit
non apparié quatre fois plus grand. Dans un régime à peu de données et à métrique
bruitée, une méthode sans entraînement achète de la \emph{puissance statistique}, et
c'est ce qui a décidé cette étude.

\subsection{Ce qui subsiste à la cardinalité matérielle}
\label{sec:res-hand}

\begin{table}[t]
\centering
\caption{Cardinalité réaliste, placement oracle. Différences appariées en bas.}
\label{tab:hand}
\begin{tabular}{lrr}
\toprule
Condition & Gain moyen & Test des signes \\
\midrule
8 ancres, position seule & $+3{,}25$ & $p=0{,}118$ \\
8 ancres $+$ normales & $+4{,}18$ & $p=0{,}118$ \\
8 ancres, placement bruité & $+1{,}78$ & $p=0{,}302$ \\
4 ancres seulement & $+0{,}70$ & $p=1{,}000$ \\
\midrule
\multicolumn{3}{l}{\emph{Différences appariées (même objet, même graine)}}\\
Normales vs sans & $+0{,}93$ & $p = 0{,}302$ \\
Bruit de placement 3\,mm / $15^\circ$ & \lo{$-2{,}40$} & \gd{$p = 0{,}035$} \\
4 zones au lieu de 8 & \lo{$-2{,}55$} & \gd{$p = 0{,}035$} \\
\bottomrule
\end{tabular}
\end{table}

Une main dextre à neuf zones n'en excite que trois à cinq dans une saisie réelle. À
une ancre par zone (tableau~\ref{tab:hand}), une réduction d'un facteur $16$ par
rapport à la campagne de référence coûte environ la moitié du plafond ($+6{,}77
\rightarrow +3{,}25$), pas la totalité. Mais les deux effets qui atteignent la
significativité sont précisément les deux propriétés d'une prise réelle : le bruit de
placement et la réduction du nombre de zones. Une vraie main a les deux, et la
condition conjointe --- ancres resserrées du simulateur, avec normales, bruit de
placement \emph{et} repère estimé plutôt qu'oracle --- n'a pas été mesurée. Nous
consignons la prédiction, avant mesure, qu'elle sera nulle.

\subsection{Résultats méthodologiques}
\label{sec:methodology}

\textbf{La variance d'une exécution à l'autre, et d'où elle vient.} Relancer huit fois
une commande identique --- même objet, même graine, mêmes ancres --- donne un
écart-type de $0{,}49$ pour la passe non ancrée mais $2{,}29$ pour la passe ancrée, et
$2{,}50$ pour la différence appariée. La lecture évidente, selon laquelle le guidage
amplifierait le non-déterminisme du GPU, est \emph{fausse}, et nous la rapportons
parce que nous l'avons d'abord tenue pour vraie. Ces répétitions ne sont pas
identiques : le repère oracle est réestimé depuis la passe non ancrée de chaque
répétition, si bien que les positions d'ancres bougent. Geler le repère sépare les
causes :

\begin{center}\small
\begin{tabular}{lcc}
\toprule
& repère réestimé & repère gelé \\
\midrule
non ancré & 0,49 & 0,68 \\
ancré & \textbf{2,29} & \textbf{0,70} \\
différence appariée & 2,50 & 1,29 \\
\bottomrule
\end{tabular}
\end{center}

Le placement figé, la passe ancrée est aussi reproductible que la passe non ancrée.
\textbf{Le guidage n'amplifie rien : 91~\% de la variance d'un score ancré vient de
la réestimation du placement des ancres.} Noter dix fois un même maillage enregistré
en explique $0{,}32$ de plus, si bien que la génération elle-même contribue pour moins
de $0{,}4$. Le problème de repère de la Sec.~\ref{sec:frame} apparaît donc deux fois
: il consomme l'essentiel du gain moyen, et il domine la variance.

Cela corrige le modèle d'erreur. Au sein d'une campagne, le repère oracle est calculé
une fois par (objet, graine) et partagé par toutes les conditions : le bruit de
placement s'annule dans les différences entre conditions ; entre campagnes il ne
s'annule pas, aucune ne partageant sa passe non ancrée. Les deux familles portent des
erreurs types de $\approx 0{,}26$ et $0{,}91$ sur 15 cas (Fig.~\ref{fig:results}a).
Nous fondons néanmoins chaque affirmation sur le test des signes : le $\sigma$ à
repère gelé repose sur huit répétitions d'un seul objet, et son intervalle de confiance
à 95~\%, $[0{,}46 ; 1{,}42]$, est assez large pour que les effets les plus petits se
déplacent entre $1{,}8$ et $5{,}5$ erreurs types selon l'endroit où se trouve la
vérité.

\textbf{Le F-Score est aveugle à la fragmentation.} Un semis d'éclats disjoints près
de la vraie surface obtient un bon score. Sur la banane : une entrée à 128 points
produit 43 composantes connexes dont la plus grosse porte 11~\% des faces, et
$47{,}83$ \Fs, contre $34{,}48$ pour une prédiction image seule qui est une surface
fermée unique. Nous rapportons le nombre de composantes et rejetons les métriques de
distance en deçà d'un seuil sur la part de la plus grosse.

\textbf{Le F-Score est aveugle aux symétries miroir.} Construire les candidats de
recalage à partir de bases issues d'une SVD donne des repères orthogonaux mais pas
nécessairement directs : \textbf{5 recalages sur 15 étaient des symétries}. Corrigé en
forçant le déterminant à $+1$.

\textbf{Un ICP mal initialisé échoue en silence.} Un ICP initialisé à l'identité
laisse le recalage hors de son bassin de convergence, sous-notant un objet de jusqu'à
$31$ points sans aucun signal d'erreur. Cela a modifié le classement entre méthodes
comparées, et c'est l'origine du verdict initial à « $-8$ points » du module appris.

% ===========================================================================
\section{Limites}
% ===========================================================================

\textbf{Le contrôle manquant.} Le résultat de visibilité dit que des ancres sur la
face visible valent autant que des ancres sur la face cachée. La face visible est
exactement ce qu'une caméra de profondeur mesure. L'expérience qui en découle ---
remplacer les contacts tactiles par une poignée de points de profondeur de la face
visible, en contraintes dans l'Éq.~\eqref{eq:loss} --- n'a pas été menée. Si elle donne
le même $+3$, la voie par contrainte n'a besoin d'aucun toucher ; si elle donne moins,
les contacts portent quelque chose de spécifique. Les deux réponses se publient, et
nous n'avons ni l'une ni l'autre.

\textbf{Un défaut hérité.} Le moteur de rendu normalise chaque objet à une dimension
maximale de $1{,}3$, alors que tous les calculs d'occultation le normalisent à $2{,}0$
en laissant la caméra en place --- l'objet est $1{,}54\times$ trop gros par rapport à
la caméra au moment où l'on décide de la visibilité. Sur le test par normale qui
étiquette les six ancres d'une prise, cela ne change \emph{rien} : aucun des 45 cas ne
bouge, et $r$ reste à $-0{,}145$. Sur la partition par suppression des points cachés,
qui sert au rappel visible/occulté, cela déplace la part visible de $0{,}7$ à $5{,}7$
points de pourcentage : ces colonnes portent cette erreur.

Les images sont des rendus Blender synthétiques, sous une seule position de caméra et
un seul éclairage. Les ancres sont échantillonnées sur le maillage de référence ---
filtrées par accessibilité, occultation et encombrement du doigt, mais sans bruit et
jamais fausses, sauf là où du bruit est ajouté délibérément. Aucun objet, capteur ou
robot réel n'intervient, bien que le banc robotique qui a motivé ce travail existe et
soit opérationnel. Une seule famille d'ossature et un seul poids de guidage ont été
utilisés ; trois graines par objet restent peu ; les expériences secondaires portent
sur cinq objets, dont l'un, la perceuse, porte une part importante de plusieurs
totaux. Les formulations 3, 4 et 7 de la taxonomie de fusion sont \emph{sans objet}
dans ce régime de données plutôt qu'invalidées --- le réentraînement à fusion précoce,
seule voie que les auteurs de l'ossature valident, était hors budget.

% ===========================================================================
\section{Conclusion}
% ===========================================================================

Un a priori génératif 3D mono-vue peut être ancré par des mesures métriques éparses,
et ce qui limite l'ancrage n'est pas les mesures. Leur nombre n'a pas d'effet sur une
plage de $16\times$ ; leurs normales n'aident pas de façon mesurable ; qu'elles soient
ou non sur la face que la caméra ne voit pas n'a pas d'effet non plus. Ce qui compte
est de savoir si elles peuvent être placées dans le repère où le générateur produit
--- et ce repère est tiré avec le bruit, si bien qu'aucun estimateur ne peut
l'apprendre. L'estimer malgré tout coûte 60 à 74~\% du gain atteignable et explique
91~\% de la variance d'une exécution à l'autre.

Fixer le repère par conditionnement est possible, et vaut en moyenne trois fois plus
que l'ancrage par contrainte, mais ce n'est pas gratuit : le modèle traite la preuve
reprojetée comme réelle, et sur un tiers des objets le prix est catastrophique, sans
sélecteur en vue qui ne connaisse déjà la réponse. Empiler des contraintes sur le
repère fixé n'ajoute presque rien.

Le mécanisme cohérent avec tout cela n'est pas que les mesures complètent ce que la
vision manque. C'est que \emph{quelques ancres vraies régularisent la sortie d'un a
priori génératif} --- position, échelle, orientation --- et que le repère propre de
l'a priori est ce qui fait obstacle. C'est une revendication plus modeste que celle
que ce travail visait, et plus générale : elle vaut pour le toucher, pour la
profondeur, et pour tout signal métrique épars qu'un robot peut apporter.

\balance
\begin{thebibliography}{99}

%% Vérifiée en ligne le 2026-09-01 : wala, ycb, smith2020, smith2021, suresh2022 et
%% dust3r ont été confrontées à arXiv / à la notice de l'éditeur. Les autres entrées
%% (dhariwal, cfg, zero123, poisson, marchingcubes, sdedit) sont des références
%% standard reproduites de mémoire -- revérifier la pagination avant soumission.

\bibitem{wala}
A.~Sanghi, A.~Khani, P.~Reddy, A.~Rampini, D.~Cheung, K.~R.~Malekshan, K.~Madan et
H.~Shayani, « Wavelet Latent Diffusion (WaLa): Billion-parameter 3D generative model
with compact wavelet encodings », \emph{arXiv:2411.08017}, 2024.

\bibitem{dhariwal}
P.~Dhariwal et A.~Nichol, « Diffusion models beat GANs on image synthesis », dans
\emph{Adv.\ Neural Inf.\ Process.\ Syst. (NeurIPS)}, 2021.

\bibitem{cfg}
J.~Ho et T.~Salimans, « Classifier-free diffusion guidance », dans \emph{NeurIPS
Workshop on Deep Generative Models}, 2021.

\bibitem{zero123}
R.~Liu, R.~Wu, B.~Van~Hoorick, P.~Tokmakov, S.~Zakharov et C.~Vondrick,
« Zero-1-to-3: Zero-shot one image to 3D object », dans \emph{Proc.\ IEEE/CVF Int.\
Conf.\ Comput.\ Vis.\ (ICCV)}, 2023.

\bibitem{dust3r}
S.~Wang, V.~Leroy, Y.~Cabon, B.~Chidlovskii et J.~Revaud, « DUSt3R: Geometric 3D
vision made easy », dans \emph{Proc.\ IEEE/CVF Conf.\ Comput.\ Vis.\ Pattern
Recognit.\ (CVPR)}, 2024.

\bibitem{ycb}
B.~Calli, A.~Singh, A.~Walsman, S.~Srinivasa, P.~Abbeel et A.~M.~Dollar, « The YCB
object and model set: Towards common benchmarks for manipulation research », dans
\emph{Proc.\ Int.\ Conf.\ Adv.\ Robot.\ (ICAR)}, 2015, p.~510--517.

\bibitem{smith2020}
E.~J.~Smith, R.~Calandra, A.~Romero, G.~Gkioxari, D.~Meger, J.~Malik et M.~Drozdzal,
« 3D shape reconstruction from vision and touch », dans \emph{Adv.\ Neural Inf.\
Process.\ Syst. (NeurIPS)}, vol.~33, 2020.

\bibitem{smith2021}
E.~J.~Smith, D.~Meger, L.~Pineda, R.~Calandra, J.~Malik, A.~Romero et M.~Drozdzal,
« Active 3D shape reconstruction from vision and touch », dans \emph{Adv.\ Neural
Inf.\ Process.\ Syst. (NeurIPS)}, 2021.

\bibitem{suresh2022}
S.~Suresh, Z.~Si, J.~G.~Mangelson, W.~Yuan et M.~Kaess, « ShapeMap 3-D: Efficient
shape mapping through dense touch and vision », dans \emph{Proc.\ IEEE Int.\ Conf.\
Robot.\ Autom.\ (ICRA)}, 2022, p.~7073--7080.

\bibitem{poisson}
M.~Kazhdan et H.~Hoppe, « Screened Poisson surface reconstruction », \emph{ACM Trans.\
Graph.}, vol.~32, n\textsuperscript{o}~3, 2013.

\bibitem{marchingcubes}
W.~E.~Lorensen et H.~E.~Cline, « Marching cubes: A high resolution 3D surface
construction algorithm », dans \emph{Proc.\ SIGGRAPH}, 1987.

\bibitem{sdedit}
C.~Meng, Y.~He, Y.~Song, J.~Song, J.~Wu, J.-Y.~Zhu et S.~Ermon, « SDEdit: Guided image
synthesis and editing with stochastic differential equations », dans \emph{Proc.\ Int.\
Conf.\ Learn.\ Represent.\ (ICLR)}, 2022.

\end{thebibliography}

\end{document}
